\begin{ex}
Khối trụ có bao nhiêu mặt phẳng đáy?
\shortans{2}
\loigiai{
Có 2 mặt đáy phẳng.
}
\end{ex}

\begin{ex}
Khối cầu có bao nhiêu đỉnh?
\shortans{0}
\loigiai{
Có 0 đỉnh.
}
\end{ex}

\begin{ex}
Khối lập phương có bao nhiêu mặt vuông?
\shortans{6}
\loigiai{
Có 6 mặt vuông.
}
\end{ex}

\begin{ex}
Khối hộp chữ nhật có bao nhiêu đỉnh?
\shortans{8}
\loigiai{
Có 8 đỉnh.
}
\end{ex}

\begin{ex}
Khối hộp chữ nhật có bao nhiêu cạnh?
\shortans{12}
\loigiai{
Có 12 cạnh.
}
\end{ex}

\begin{ex}
Trong các khối: khối lập phương, khối hộp chữ nhật, khối trụ, khối cầu. Có bao nhiêu khối có đỉnh nhọn?
\shortans{2}
\loigiai{
Khối lập phương và khối hộp chữ nhật có đỉnh nhọn (2 khối).
}
\end{ex}

\begin{ex}
Trong các khối trên, có bao nhiêu khối có thể lăn được trên mặt bàn phẳng?
\shortans{2}
\loigiai{
Khối trụ và khối cầu có thể lăn được (2 khối).
}
\end{ex}

\begin{ex}
Trên giá đồ chơi có: 4 khối cầu, 3 khối trụ, 2 khối lập phương. Có tất cả bao nhiêu khối hình?
\shortans{9}
\loigiai{
Có tất cả: $4 + 3 + 2 = 9$ khối hình.
}
\end{ex}

\begin{ex}
Số khối cầu nhiều hơn số khối lập phương bao nhiêu khối?
\shortans{2}
\loigiai{
Nhiều hơn: $4 - 2 = 2$ khối.
}
\end{ex}

\begin{ex}
Xếp 3 khối lập phương chồng lên nhau được một cột cao. Xếp 2 khối trụ chồng lên nhau được một cột khác. Có tất cả bao nhiêu khối hình được dùng?
\shortans{5}
\loigiai{
Có: $3 + 2 = 5$ khối.
}
\end{ex}

\begin{ex}
Một người thợ tiện gỗ làm được 10 quả bóng gỗ dạng khối cầu. Bác chia đều vào 2 hộp. Mỗi hộp có bao nhiêu quả bóng?
\shortans{5}
\loigiai{
Mỗi hộp có: $10 : 2 = 5$ quả bóng.
}
\end{ex}

\begin{ex}
Một nhà máy sản xuất lon nước ngọt đóng thành từng lốc, mỗi lốc 6 lon hình khối trụ. Hỏi 3 lốc như thế có bao nhiêu lon nước ngọt?
\shortans{18}
\loigiai{
Có: $6 \times 3 = 18$ lon nước ngọt.
}
\end{ex}

\begin{ex}
Một hộp đựng bóng tennis dạng khối trụ chứa vừa vặn 3 quả bóng tennis dạng khối cầu. Hỏi 4 hộp như thế chứa được bao nhiêu quả bóng tennis?
\shortans{12}
\loigiai{
Chứa được: $3 \times 4 = 12$ quả bóng.
}
\end{ex}

\begin{ex}
Khối trụ có mấy mặt cong xung quanh?
\shortans{1}
\loigiai{
Khối trụ có 1 mặt cong xung quanh.
}
\end{ex}

\begin{ex}
Khối cầu có mấy mặt cong bao quanh?
\shortans{1}
\loigiai{
Khối cầu có 1 mặt cong khép kín bao quanh.
}
\end{ex}

\begin{ex}
Tổng số mặt đáy phẳng của 4 khối trụ là bao nhiêu?
\shortans{8}
\loigiai{
Có: $2 \times 4 = 8$ mặt đáy.
}
\end{ex}

\begin{ex}
Tổng số đỉnh của 2 khối hộp chữ nhật là bao nhiêu?
\shortans{16}
\loigiai{
Có: $8 + 8 = 16$ đỉnh.
}
\end{ex}

\begin{ex}
Tổng số mặt của 2 khối lập phương là bao nhiêu?
\shortans{12}
\loigiai{
Có: $6 + 6 = 12$ mặt.
}
\end{ex}

\begin{ex}
Một cây nến hình khối trụ có chiều cao 15 cm. Sau khi thắp sáng một lúc, cây nến bị cháy ngắn bớt 6 cm. Cây nến còn lại dài bao nhiêu xăng-ti-mét?
\shortans{9}
\loigiai{
Còn lại: $15 - 6 = 9$ cm.
}
\end{ex}

\begin{ex}
Bác thợ mộc có một khúc gỗ dài 20 dm dạng khối trụ. Bác cưa thành 4 khúc bằng nhau. Mỗi khúc dài bao nhiêu đề-xi-mét?
\shortans{5}
\loigiai{
Mỗi khúc dài: $20 : 4 = 5$ dm.
}
\end{ex}

\begin{ex}
Bạn Hùng có 14 viên bi ve khối cầu, bạn Dũng có ít hơn bạn Hùng 5 viên. Bạn Dũng có bao nhiêu viên bi ve?
\shortans{9}
\loigiai{
Dũng có: $14 - 5 = 9$ viên bi.
}
\end{ex}

\begin{ex}
Có bao nhiêu khối cầu trong một hộp chứa 8 quả bóng bàn và 6 chiếc vợt cầu lông?
\shortans{8}
\loigiai{
Có 8 quả bóng bàn dạng khối cầu.
}
\end{ex}

\begin{ex}
Xếp 4 khối trụ thành một hàng ngang, cần thêm bao nhiêu khối trụ nữa để có 10 khối trụ?
\shortans{6}
\loigiai{
Cần thêm: $10 - 4 = 6$ khối.
}
\end{ex}

\begin{ex}
Một chiếc bánh sinh nhật dạng khối trụ tròn có 2 tầng bằng nhau. Cả hai tầng có bao nhiêu mặt đáy phẳng tròn ngoài cùng?
\shortans{2}
\loigiai{
Có 2 mặt đáy phẳng ngoài cùng.
}
\end{ex}

\begin{ex}
Số cạnh của một khối cầu là bao nhiêu?
\shortans{0}
\loigiai{
Khối cầu có 0 cạnh.
}
\end{ex}
